\chapter{Idiome und Redewendungen}
\label{cha:Idiome}

Dieses Kapitel enthält häufige deutsche Redewendungen und idiomatische Ausdrücke.

\section{Alltagsidiome}
\label{sec:Alltagsidiome}

\bigskip
\begin{center}
\begin{tabular}{|l|p{8cm}|}
\hline
Idiom & Bedeutung \\ \hline
Das ist mir schnuppe & Das ist mir egal. \\ \hline
Das ist mir scheißegal & Das ist mir völlig egal (vulgär). \\ \hline
In die Röhre gucken & Keinen Erfolg haben, enttäuscht werden. \\ \hline
Jemanden auf dem Kieker haben & Jemanden ständig kritisieren. \\ \hline
Den Bock zum Gärtner machen & Die falsche Person für eine Aufgabe einsetzen. \\ \hline
Das Handtuch werfen & Aufgeben, aufhören. \\ \hline
Jemanden über den Tisch ziehen & Jemanden betrügen. \\ \hline
Etwas an der Backe haben & Schulden haben. \\ \hline
Den Nagel auf den Kopf treffen & Genau das Richtige sagen. \\ \hline
Aus der Haut fahren & Wütend werden. \\ \hline
Die Krise kriegen & Panik bekommen. \\ \hline
Sich etwas unter den Nagel reißen & Etwas stehlen. \\ \hline
Den Bogen raus haben & Etwas gut können. \\ \hline
Das Rad nicht neu erfinden & Keine neue Lösung suchen, wenn es schon eine gibt. \\ \hline
Jemanden in die Pfanne hauen & Jemanden verraten. \\ \hline
Etwas auf die hohe Kante legen & Etwas sparen. \\ \hline
Den Spieß umdrehen & Die Situation umkehren. \\ \hline
Sich etwas durch den Kopf gehen lassen & Etwas überlegen. \\ \hline
Die Ärmel hochkrempeln & Sich an die Arbeit machen. \\ \hline
Ins kalte Wasser springen & Etwas ohne Vorbereitung beginnen. \\ \hline
\end{tabular}
\end{center}

\section{Idiome mit Körperteilen}
\label{sec Koerper}

\bigskip
\begin{center}
\begin{tabular}{|l|p{8cm}|}
\hline
Idiom & Bedeutung \\ \hline
Den Kopf verlieren & In Panik geraten. \\ \hline
Den Kopf zerbrechen & Sich den ganzen Tag mit etwas beschäftigen. \\ \hline
Das Herz auf der Zunge tragen & Offen und ehrlich sein. \\ \hline
Das Herz in der Hand haben & Sensibel und gefühlvoll sein. \\ \hline
Aus dem Bauch heraus & Intuitiv, ohne nachzudenken. \\ \hline
Bauchschmerzen haben & Sich Sorgen machen. \\ \hline
Sich etwas durch den Kopf schlagen & Über etwas nachdenken. \\ \hline
Den Finger auf die Wunde legen & Das Problem direkt ansprechen. \\ \hline
Die Nase voll haben & Genervt sein von etwas. \\ \hline
Die Nase hoch tragen & Überheblich sein. \\ \hline
Mit der Nase auf etwas stoßen & Jemanden auf ein Problem aufmerksam machen. \\ \hline
Sich etwas aus den Fingern saugen & Sich etwas ausdenken. \\ \hline
Daumen drücken & Glück wünschen. \\ \hline
Die Daumen drehen & Faulenzen. \\ \hline
Die Stirn bieten & Mutig Widerstand leisten. \\ \hline
Ins Gesicht sagen & Offen die Wahrheit sagen. \\ \hline
Das Gesicht wahren & Den Schein wahren. \\ \hline
Sich etw. nicht aus der Ruhe bringen lassen & Gelassen bleiben. \\ \hline
Ein地方 auf dem Kopf haben & Verrückt sein. \\ \hline
Sich etwas hinter die Ohren schreiben & Sich etwas merken. \\ \hline
\end{tabular}
\end{center}

\section{Idiome mit Tieren}
\label{sec:Tiere}

\bigskip
\begin{center}
\begin{tabular}{|l|p{8cm}|}
\hline
Idiom & Bedeutung \\ \hline
Der Fuchs & Schlauer, listiger Mensch. \\ \hline
Das schwarze Schaf & Außenseiter, Sündenbock. \\ \hline
Der Vogel ist ausgeflogen & Die Person ist weg. \\ \hline
Ein alter Hase & Erfahrene Person. \\ \hline
Da liegt der Hund begraben & Das ist das eigentliche Problem. \\ \hline
Jemanden wie einen Hund behandeln & Schlecht behandeln. \\ \hline
Der Hund in der Mühe sein & Ein großes Problem sein. \\ \hline
Keinen Hund mehr raus lassen & Total erschöpft sein. \\ \hline
Wie ein Hund leben & Ein schweres Leben führen. \\ \hline
Die Katze aus dem Sack lassen & Ein Geheimnis verraten. \\ \hline
Das ist Katz und Maus & Das ist ein kompliziertes Spiel. \\ \hline
Die Katze hat die Zunge verschluckt & Sprachlos sein. \\ \hline
Der Hahn im Korb sein & Der einzige Mann unter Frauen. \\ \hline
Das ist ein anderes Kaliber & Das ist etwas ganz anderes. \\ \hline
Da steppt der Bär & Da wird gefeiert. \\ \hline
Einen Bärenhunger haben & Sehr hungrig sein. \\ \hline
Den Bogen raus haben & Etwas gut können. \\ \hline
Auf den Bäumen klettern können & Vielseitig begabt sein. \\ \hline
Die Flöhe husten hören & Überempfindlich sein. \\ \hline
Wütend wie ein Stier & Sehr wütend. \\ \hline
\end{tabular}
\end{center}

\section{Feste Verbindungen}
\label{sec:FesteVerbindungen}

\bigskip
\begin{center}
\begin{tabular}{|l|p{8cm}|}
\hline
Wendung & Bedeutung \\ \hline
Das Gegenteil von & Das Entgegengesetzte. \\ \hline
Mit anderen Worten & Anders ausgedrückt. \\ \hline
Auf den ersten Blick & Sofort, beim ersten Anblick. \\ \hline
In der Tat & Tatsächlich, wirklich. \\ \hline
Das sei dahingestellt & Das bleibt offen. \\ \hline
Alles in allem & Zusammengefasst. \\ \hline
Im Großen und Ganzen & Im Allgemeinen. \\ \hline
Ohne Weiteres & Ohne Probleme. \\ \hline
Mit Recht & Berechtigt, zu Recht. \\ \hline
Im Prinzip & Grundsätzlich. \\ \hline
Im Allgemeinen & Normalerweise. \\ \hline
Vor allem & Besonders, hauptsächlich. \\ \hline
Zumindest & Mindestens. \\ \hline
Ganz zu schweigen von & Nicht zu vergessen. \\ \hline
Geschweige denn & Erst recht nicht. \\ \hline
Nichtsdestotrotz & Trotzdem. \\ \hline
Zugutekommen & Von Vorteil sein. \\ \hline
Ins Gewicht fallen & Wichtig sein. \\ \hline
Auf dem Spiel stehen & In Gefahr sein. \\ \hline
In die Bresche springen & Helfen, einspringen. \\ \hline
\end{tabular}
\end{center}

\bigskip
\noindent\textbf{Grußformeln und Höflichkeiten:}

\begin{itemize}
    \item \textbf{Mahlzeit!} - Guten Appetit (auch als Gruß mittags)
    \item \textbf{Prost!} - Zum Wohl (bei Alkohol)
    \item \textbf{Guten Rutsch!} - Guten Start ins neue Jahr
    \item \textbf{Toi, toi, toi!} - Viel Glück!
    \item \textbf{To break a leg!} - Viel Erfolg!
    \item \textbf{Pfiat di!} - Servus (bayerisch)
    \item \textbf{Grüß Gott!} - Guten Tag (süddeutsch)
    \item \textbf{Moinsen!} - Moin (norddeutsch)
\end{itemize}

\chapterend

\chapter{Falsche Freunde (False Friends)}
\label{cha:FalseFriends}

Dieses Kapitel listet Wörter, die ähnlich aussehen wie englische Wörter, aber eine andere Bedeutung haben.

\section{Gefährliche Falschfreunde}
\label{sec:GefaesslicheFF}

\bigskip
\begin{center}
\begin{tabular}{|l|p{5cm}|p{5cm}|}
\hline
Deutsch & Englisch (falsch) & Englisch (richtig) \\ \hline
aktuell & actual & current, present \\ \hline
Billion & billion & trillion \\ \hline
eventuell & eventful & possibly, perhaps \\ \hline
fabelhaft & fable-like & fabulous, great \\ \hline
Gift & gift & poison \\ \hline
klug & clown & clever, smart \\ \hline
Knabe & knave & boy, lad \\ \hline
Mega & mega & great, huge (slang) \\ \hline
Pelz & pelt & fur, coat \\ \hline
Rotz & rotz & snot, mucus \\ \hline
Schaden & shadow & damage, harm \\ \hline
See & sea & lake \\ \hline
Stadion & stadium & (same in German) \\ \hline
Tonne & ton & barrel, bin \\ \hline
wahrhaft & worried & truly, indeed \\ \hline
\end{tabular}
\end{center}

\section{Häufige Falschfreunde}
\label{sec:HaeufigFF}

\bigskip
\begin{center}
\begin{tabular}{|l|p{5cm}|p{5cm}|}
\hline
Deutsch & Englisch (falsch) & Englisch (richtig) \\ \hline
actually & actually & eigentlich, tatsächlich \\ \hline
angeblich & allegedly & supposedly, supposedly \\ \hline
Bibliothek & bibliotheca & library \\ \blich{blau} & blue & \\ \hline
Fabrik & fabric & factory \\ \hline
fast & fast & almost, nearly \\ \hline
gebrauchen & grab/use & to use \\ \hline
Gymnasium & gymnasium & secondary school \\ \hline
hell & hell & light, bright \\ \hline
Hose & hose & pants, trousers \\ \hline
Karosse & carosse & carriage, coach \\ \hline
Konzept & concept & draft, plan \\ \hline
Konto & conto & account \\ \hline
Kurve & curve & bend, turn \\ \hline
letzt & last & last, final \\ \hline
Messe & mass & trade fair, mass (church) \\ \hline
Miete & mite & rent \\ \hline
Milliarde & milliard & billion \\ \hline
Nummer & number & number, size \\ \hline
Pfund & pound & pound (weight/currency) \\ \hline
Platz & place & square, seat, space \\ \hline
praktisch & practical & practical, handy \\ \hline
Rat & rat & advice, council \\ \hline
real & real & realistic, real \\ \hline
Ring & ring & ring, arena \\ \hline
riskant & risky & risky, dangerous \\ \hline
Rohr & row & pipe, tube \\ \hline
Scheck & check & cheque, check \\ \hline
sensibel & sensible & sensitive \\ \hline
Sympathie & sympathy & liking, affinity \\ \hline
Tempo & tempo & speed, pace \\ \hline
\end{tabular}
\end{center}

\section{Ausdrücke und Phrasen}
\label{sec:FFAusdruecke}

\bigskip
\begin{center}
\begin{tabular}{|l|p{6cm}|p{6cm}|}
\hline
Deutsch & Bedeutung (falsch) & Bedeutung (richtig) \\ \hline
Ich bin auf dem Weg & I am on the way & I'm on my way \\ \hline
Das ist sehr sensible & That is very sensible & That is very sensitive \\ \hline
Ich habe keine Ahnung & I have no idea & I have no idea (correct!) \\ \hline
Das ist die Situation & That is the situation & That's the situation (correct) \\ \hline
Ich bin einverstanden & I am agreed & I agree / I'm in agreement \\ \hline
Lass uns & Let us & Let's (imperative) \\ \hline
Was machst du? & What do you make? & What are you doing? \\ \hline
Das ist alles & That is all (all = alles) & That's all (correct) \\ \hline
Allerdings & All the ditches & However, indeed \\ \hline
Gewöhnlich & Usually (gewöhnlich) & ordinarily, usually \\ \hline
Zunächst & At first (zunächst) & first, initially \\ \hline
Ziemlich & Tightly & quite, rather, pretty \\ \hline
\end{tabular}
\end{center}

\chapterend